\documentclass[11pt,a4paper]{article}
\usepackage[T1]{fontenc}
\usepackage{lmodern}
\usepackage[margin=22mm]{geometry}
\usepackage{amsmath,amssymb,amsthm,microtype}
\usepackage[hidelinks]{hyperref}
\setlength{\parindent}{0pt}
\setlength{\parskip}{5pt}
\setlength{\emergencystretch}{2em}
\newtheorem{lemma}{Lemma}
\title{The Nine-Dimensional Snake Problem:\\A Short Explanation of the Current Argument}
\author{Iori Hironaka \quad Tommy Ma \quad Steven Zhan}
\date{Research state after round 209}
\begin{document}
\maketitle

\textbf{Problem and present result.}
A snake in $Q_9$ is an induced path: two selected vertices are adjacent
exactly when they are consecutive on the path. Its length $m$ counts
edges. A verified construction has 191 edges. The project retains an
archived contributor proof of the upper bound 235, whose complete
independent rebuild remains pending. The maximum is unresolved.
The numerical prediction 196 has neither a construction nor an optimality
proof. The current construction investigation compares dimensions six
through nine, with general restrictions on two-parent constructions
and no new witness. A whole-path classification now restricts changes
to one or two coordinate histories. A new native enumeration excludes
a different class of the saved 98-edge parent even with five arbitrary
coordinate histories free. Its precise scope and proof-assistant
boundaries are explained in Section 6. The current sufficient pending-share profile is $(15,14,14,14)$
for $012034,012301,012340,012345$, respectively. Together with the retained
local and graph premises, it implies $m\le234$ without a positive-loss
assumption. All four required numerical floors have native contradiction
certificates; their complete Lean and geometric interpretation, and the
inherited graph requirements, remain unfinished. A separate contributor chain gives conditional 233; the new boundary
argument weakens its local share requirements. Its full certificate
assembly and the new geometric integration remain pending. The argument
and these limitations are given in order.

\textbf{1. The exact square budget.}
Assume $m\ge1$; the single-vertex path satisfies the length bounds directly.
Write $v=m+1$, let $T$ count coordinate squares with exactly one selected
corner, and let $s_0$ count empty squares. There are
$\binom92 2^7=4608$ coordinate squares, and each vertex lies in 36 of them.
A snake contains no fully selected square. Its three-selected-corner
squares correspond exactly to its $v-2$ internal vertices. Counting
squares and selected corner incidences therefore gives
\[
\boxed{T+2s_0+35(m+1)=9214.}
\]
This identity converts a lower count estimate into an upper length estimate.

Let $\alpha_i$ be the singleton-square count at the $i$th path vertex.
Every count is positive, the endpoint counts are at least four, an internal
count one has internal neighbours at least four, and three consecutive
internal counts cannot all equal two. These are established geometric
results. They imply that six internal counts sum to at least fourteen and
seven sum to at least sixteen.

Here is the short deduction. Every adjacent pair sums to at least four.
A triple sums to at least six, with equality only at $(1,4,1)$; four
counts sum to at least nine, since two pairs each summing to four would
force four twos. Split six counts into two triples. If one has sum six,
its adjacent count in the other triple is at least four and the remaining
pair contributes at least four. Otherwise both triples sum to at least
seven. For seven counts, split into four and three. The ordinary case
contributes at least $9+7$; a final $(1,4,1)$ forces at least $6+4+6$.
These deductions and their application to actual snake prefixes pass Lean.

The archived count-only potential gives
$4(m-7)\le\sum_{j=7}^{m-1}\min\{6,\alpha_j\}+7$ for $m\ge7$.
The six omitted internal counts and the endpoints now contribute at
least $14+8=22$, yielding $T\ge4m-13$. Shorter nontrivial paths satisfy
this estimate directly. Substitution into the square identity gives
\[
\boxed{39m+2s_0\le9192.}
\]
Thus $m\le235$, and $m=235$ requires $s_0\le13$. Proving fourteen empty
squares in that case would suffice to reach 234; the earlier target was
eighteen. This deduction uses the archived walk estimate. Its new boundary
and arithmetic are checked; the archive's full independent rebuild is
still pending. Seven disjoint share-96 windows, or 33 such windows with
distinct starts allowing overlap, would also force fourteen empty squares.
Their required occurrence frequencies remain unproved.

\textbf{2. Retaining the empty-square contribution exactly.}
For a cube vertex $q$, let $e(q)$ count incident empty squares and $d(q)$
count selected neighbours. Each selected vertex $x$ receives the integer share
\[
\sigma_{12}(x)=\sum_{q\sim x}\left\lfloor\frac{12e(q)}{d(q)}\right\rfloor.
\]
Every denominator here is positive because $x$ is selected. Define the
remainder at $q$ to be $12e(q)\bmod d(q)$ for $d(q)>0$, and $12e(q)$
otherwise; let $L$ be their sum. Exactly $d(q)$ selected vertices receive
$q$'s share, and the division algorithm recovers $12e(q)$ after the
remainder is added. Each empty square has four corners, so
\[
\sum_{x\in S}\sigma_{12}(x)+L=48s_0.
\]
Writing $G=24T+\sum_{x\in S}\sigma_{12}(x)$ and multiplying the square
identity by 24 gives
\[
\boxed{G+840(m+1)+L=221136.}
\]
The physical counts, exact share-plus-loss identity and sufficient budget
inequality are now connected in one checked Lean proof.

\textbf{3. From local windows to a global length bound.}
The nine-transition model charges $120\min\{6,\alpha_{r+8}\}$ and a
valid lower bound on the five shares at positions $r+4,\ldots,r+8$.
Assume $m\ge8$; shorter paths already satisfy the stated upper bounds.
Let $W$ be the sum over all $m-8$ windows. The omitted singleton positions
are $0,1,\ldots,7,m$. Their counts sum to at least $16+8=24$.
Each share appears in at most five windows. Consequently,
\[
W+2880\le5G,\qquad W+4200m+5L\le1098600.
\]
The earlier proof used only 21 omitted counts. The new estimate is checked
for actual snakes, with each local share lower bound retained as a premise.

For a nonnegative mean $\mu$, suppose every traced edge satisfies
$w(u,v)+p(u)-p(v)\ge\mu$ and the potential range is at most $R\ge0$,
summing cancels intermediate potential values and yields
$\mu(m-8)\le W+R$. Combining the inequalities gives
\[
\boxed{(4200+\mu)m+5L\le1098600+8\mu+R.}
\]
The old potential $\mu=4462/9$, $R=5411/9$ therefore already implies
$m\le235-3063/42262<235$ using only $L\ge0$.
A separate proof of $L\ge4$ is no longer needed for this conditional route.

\textbf{4. One critical floor fifteen gives the conditional bound.}
A direction pattern records successive coordinate flips, with labels renamed
consistently. A share floor bounds the sum at the associated five window
centres. The finite retained graph has 4,661,572 edges. The current sufficient
floors, ordered as $012034,012301,012340,012345$, are $(15,14,14,14)$.
The graph construction retains all its previously stated count and exclusion
premises.

The weaker pending-share profile
\[
(15,14,14,14),\qquad\text{for }012034,012301,012340,012345,
\]
admits a retained-graph potential with $\mu=4454/9$ and $R=5421/9$.
All 4,661,572 edge inequalities passed an exact integer scan, after
reconstruction reproduced three archived weight hashes. Substituting
in the boundary-24 comparison gives
\[
\boxed{42254m+45L\le9928453.}
\]
Since $L\ge0$ and $235\cdot42254=9928453+1237$, this yields
\[
m\le235-\frac{1237}{42254}<235,\qquad m\le234.
\]
Thus only the critical pending floor must reach fifteen; the other three
may remain fourteen. No separate exclusion of $L=0$ or $L=1$ is required
for this sufficient profile. Lean also checks the implication for actual snakes, with the local-window
and potential inequalities retained as explicit premises.

A reconstructed nine-edge cycle has weight 4454, and a 24-edge walk has
weight 11275. Hence every admissible mean and range for this fixed graph
satisfies $9\mu\le4454$ and $R\ge24\mu-11275$. Consequently,
\[
\frac{1098600+8\mu+R}{4200+\mu}
\ge\frac{1087325+32\mu}{4200+\mu}
\ge\frac{9928453}{42254}.
\]
The middle expression decreases with $\mu$, and the new potential attains
the final value at $\mu=4454/9$. The denominator-cleared optimality
deduction passes Lean. This proves optimality of the loss-free comparison
for the fixed profile, boundary and zero occurrence credit.

Fifteen is also the least sufficient integer floor for $012034$ when the
other pending floors stay fourteen in this argument. At uniform floor
fourteen, the saved cycle and walk force $9\mu\le4448$ and
$R\ge8\mu-3356$, keeping the loss-zero comparison strictly above 235.
The strict arithmetic obstruction passes Lean; lowering that floor further
only weakens edge weights.

Every numerical floor in the new profile has a complete native Boolean
refutation. The $012301$ floor-fourteen input uses the necessary extension
condition for snakes with at least seven edges; shorter snakes already
satisfy the desired global bound. The full Lean contradiction chains,
geometric interpretation and inherited count and graph premises remain
unfinished. The optimized comparison therefore remains conditional.

An earlier sufficient profile raised all four pending floors to fifteen.
The new profile weakens three of those requirements. Its purpose is to
reduce the local proof obligations while retaining the conditional bound 234.

\textbf{5. An alternative route using only floor fourteen.}
If the critical floor is left at fourteen, the following earlier
comparisons still apply under their retained graph and count premises:
\[
\begin{array}{c|c|c|c}
\text{floors}&9\mu&9R&\text{inequality}\\\hline
(14,16,16,16)&4448&5346&42248m+45L\le9928330\\
(14,14,14,14)&4448&5380&42248m+45L\le9928364
\end{array}
\]
Every inequality on the 4,661,572 retained graph edges passed an exact
integer scan. At $m=235$, even the weaker second row forces
$45L\le84$, hence $L\le1$. Thus uniform floor fourteen and $L\ge2$
suffice for the conditional upper bound 234. The remaining hypothetical
case has exactly 235 edges and $L=0$ or $L=1$.

Two is the least sufficient integer loss threshold for these fixed
comparisons. A nine-edge cycle of weight 4448 forces $9\mu\le4448$.
A nine-edge walk of weight 3854 in the first profile, and an eight-edge
walk of weight 3356 in the second, force the respective range bounds.
With $L=1$, the length comparison is
\[
\frac{1098595+8\mu+R}{4200+\mu}.
\]
Substituting $R\ge9\mu-3854$ or $R\ge8\mu-3356$ gives, respectively,
\[
\frac{1094741+17\mu}{4200+\mu},\qquad
\frac{1095239+16\mu}{4200+\mu}.
\]
Both decrease as $\mu$ increases, so the cycle bound makes their minima
occur at $\mu=4448/9$. The resulting values are
$235+5/42248$ and $235+39/42248$. Both exceed 235, and the new potentials
attain these minima. Thus changing the mean and range alone cannot make
one unit of loss sufficient for these fixed graph and boundary comparisons. Lean checks
these deductions and their actual-snake implications under explicit window
and potential premises; scalar reconstruction checks the witness rows.

Integer rounding of the aggregate constraints also leaves the case open:
\[
(m,L,G,W)=(235,1,22895,111595)
\]
satisfies the global budget, packing constraint and both potential
inequalities. This arithmetic witness is Lean-checked. It is not a
snake construction; it identifies a limitation of those aggregate
constraints alone.

Both remaining loss cases dominate the cube, meaning every cube vertex
has a selected neighbour. At $L=0$ every corner
remainder vanishes. At $L=1$, exactly one is positive; it must have
five selected neighbours and three incident empty squares. Indeed,
$e\le\binom{9-d}{2}$, and the finite possibilities with
$12e\bmod d=1$ give only $(d,e)=(5,3)$. This arithmetic is Lean-checked.
Neither remaining case has been excluded.

For a set $C$ of selected window centres, let $k(q)$ count the selected
neighbours of $q$ belonging to $C$, and let $r(q)$ be its remainder.
Since $0\le k(q)\le d(q)$, double counting gives
\[
\sum_{x\in C}\sigma_{12}(x)
=\sum_{q:d(q)>0}\frac{12e(q)k(q)}{d(q)}
-\sum_{q:d(q)>0}\frac{r(q)k(q)}{d(q)},
\qquad
0\le\sum_{q:d(q)>0}\frac{r(q)k(q)}{d(q)}\le L.
\]
This bounds the local rounding deficit by $L$. The sharper incidence
argument below reduces the sufficient rational-share target to a
strict bound greater than $72/5$ when $L\le1$. The double-counting
deduction is written here; the required rational lower bound remains unproved.

\paragraph{Sharper rounding bound for the pending patterns.}
For each of $012034,012301,012340,012345$, a cube vertex is adjacent to
at most two of the five window centres. Three neighbours would have to
be the first, third and fifth centres, since all neighbours have the same
parity. The first and fifth centres differ in four coordinates in these
patterns and therefore cannot share a neighbour. Canonical finite
neighbour enumerations also pass Lean; transfer by translation and
coordinate permutation is the written geometric argument.

When $L=1$, the only positive remainder has degree five. Its contribution
to a window's rational-to-integer deficit is therefore at most $2/5$.
Thus a strict rational share bound greater than $72/5$ suffices for
integer share at least fifteen; at $L=0$, a strict bound greater than
fourteen suffices. Equality at $72/5$ does not suffice by this argument.
Clearing denominators by 2520 gives $A=2520S+B$, with $B\le1008$,
so $A>36288$ implies $S\ge15$. These arithmetic statements pass Lean.
The required rational lower bound remains an open local proof target.

\paragraph{Path restrictions for the remaining window.}
The radius-four share-cap-thirteen queries for $012034,012340,012345$
now have complete native proof replays. Every added clause and the final
contradiction were checked against the saved input. This verification
does not yet supply a Lean contradiction proof or the full geometric
interpretation of the encoding.

The original $012301$ relaxation admits a checked selected set with all
five shares zero and zero loss over the window-neighbour set. It has
sixteen induced eight-cycle components and total loss 5922. Hence a
local remainder budget alone does not remove this obstruction; the
configuration does not satisfy the hypotheses of an actual low-loss snake.

In canonical coordinates, the window is $0,1,3,7,15,14,12$. Its endpoints
have common neighbours 4 and 8. A common neighbour of two path vertices
separated by more than two edges cannot belong to an induced path:
its path index would have to differ by one from both indices, forcing
their separation to be at most two. Thus both closing vertices are absent.
Moreover, a six-edge window in a snake of length at least seven extends
at one end. If the window starts after the first path vertex, use its
predecessor; otherwise use the successor of its last vertex. These
two path properties and the canonical exclusions pass Lean checks.

The closing exclusions alone leave a checked assignment with integer
share 6, rational share $52/7$ and window-neighbour loss 5. In that
assignment, the entire window is an isolated seven-vertex component.
Adding the necessary extension condition removes this example. The
resulting share-cap-thirteen query returns a native contradiction.
The complete saved proof now passes native replay, including its final
empty clause. All four floor-fourteen inputs therefore have complete
native refutation checks.
The extension condition is valid for every actual snake of length at
least seven, so it covers all lengths relevant to the upper-bound argument.
The full Lean replay and the transfer from arbitrary coordinates to the
Boolean model remain pending. No unconditional numerical bound follows
from these local computations alone.

\paragraph{The extension uses an unused coordinate.}
The possible extension can be restricted further. At the endpoint 0,
coordinates 0 and 1 give the existing neighbour 1 or the forbidden vertex 2.
Selecting 2 would give vertex 3 three selected neighbours: 1, 7 and 2.
Coordinates 2 and 3 give the excluded closing vertices 4 and 8.
At endpoint 12, the first four coordinates give 13, 14, 8 and 4.
Here 14 is the existing neighbour, and selecting 13 would give vertex 15
the three selected neighbours 7, 14 and 13. Therefore every external
extension uses a coordinate $a$ with $4\le a<9$.
The degree exclusions and this canonical-window statement pass Lean checks.

Permuting the five unused coordinates and, if necessary, reversing the
window reduces the extension to vertex 16 at the first end. The reversal
is normalized by translation by 12 followed by the coordinate swaps
$0\leftrightarrow1$ and $2\leftrightarrow3$. This preserves the five-centre
share sum. The representative extended segment is
$16,0,1,3,7,15,14,12$, with relabelled transition word $0123412$.
The full isometry transfer remains a written argument.

The floor-fifteen query allowing either extension reached its proof-size
limit and gives no conclusion. With the representative extension fixed,
the query excluding share at most fourteen returned a contradiction.
Its complete native proof replay passed. This is stronger local Boolean
evidence for $012301$; its full Lean and geometric interpretation remain
unfinished.

\paragraph{Why this single improvement does not reduce the comparison.}
In the retained profile $(14,15,14,14)$, ordered as
$012034,012301,012340,012345$, the saved nine-edge cycle still has weight
4448: none of its edges has shape $012301$. The saved eight-edge walk
has two such edges, so its weight rises from 3356 to 3358. Direct
reconstruction of the source rows checks both identities. Every mean
and range for this fixed graph must therefore satisfy
\[
9\mu\le4448,\qquad R\ge8\mu-3358.
\]
For $L=1$, these constraints imply
\[
\frac{1098595+8\mu+R}{4200+\mu}
\ge235+\frac{21}{42248}>235.
\]
Lean checks this comparison for arbitrary nonnegative rational parameters
after clearing denominators. The $L=0$ comparison is larger still.
Thus raising only the $012301$ floor to fifteen cannot yield 234 through
the current mean-and-range argument in the two remaining loss cases.
This is an obstruction from saved walks, with no claim that a new
potential attains the displayed lower comparison.

\paragraph{A smaller certificate for the critical floor fifteen.}
The radius-four query for $012034$ excluding integer share at most fourteen
now has a complete native refutation. It uses the broad count intervals
$[0,21]$ and the existing local constraints, with no additional extension
or cycle clauses. Its complete replay checks 119,146 deductions.
Extracting the dependency closure of the final contradiction reduces this
to 12,424 original clauses and 52,489 deductions, with 3,619,677 propagation
hints. Both this core and its dense variable renaming pass separate full
native replays. The dense proof has 5,342 variables; a last-use schedule
needs at most 12,593 stored clauses. These sizes fit the existing Lean
checker's data format, while the actual kernel replay and geometric
interpretation remain unfinished.

This certificate addresses the shape occurring six times on the critical
nine-edge cycle. It supplies the critical local input for the weaker profile $(15,14,14,14)$,
in the order $012034,012301,012340,012345$. The full retained-graph calculation in Section 4 now establishes its
sufficiency under the stated premises. The local certificate alone
does not establish the global comparison.

\paragraph{Reordering threshold terms in the critical certificate.}
The streamed $012034$ input appends the fifteenth threshold bit for each
of its 35 corner occurrences after the first fourteen layers. The earlier
clause checker required the threshold terms in a fixed order. Its
soundness proof now allows any permutation of those terms, preserving
multiplicities. Indeed, for nonnegative integer corner shares $s_j$,
\[
\sum_{j=1}^{35}\sum_{t=1}^{15}\mathbf{1}_{\{t\le s_j\}}
 =\sum_{j=1}^{35}\min(s_j,15)\le\sum_{j=1}^{35}s_j.
\]
A permutation preserves the left-hand count. Thus a physical window
share at most fourteen permits an assignment satisfying the cap clause,
which negates the counter output for at least fifteen true terms.
The internal counter implications follow from
$i\le A$ and $j\le B$ implying $i+j\le A+B$ for the two child counts.
The selected-vertex, endpoint and empty-square assumptions are unchanged.

The generalized counting lemma and clause-soundness theorem pass Lean.
The reconstructed input also matches all 12,424 clauses of the saved
dense core exactly. At the end of round 185, its concrete interpretation remained
pending: the counter and totalizer structure checks exceeded the existing
256 MiB limit when checked as single propositions. Splitting constant
data and supplying explicit meaning trees has not removed that
obstruction. No concrete clause batch or final contradiction was
accepted in round 185. The next step was to prove those structural
conditions in smaller parts, then check the 98 clause batches and their
composition. The numerical bound 234 retains its stated assumptions.

\paragraph{Dividing the concrete interpretation into smaller proofs.}
The endpoint counter contains 286 distinct vertex flags. Its literal
interpretation is now checked in eighteen blocks. Strict ordering is
proved within each block and between each block and its suffix; a general
append lemma combines these results. This establishes the entire endpoint
counter condition while retaining the existing memory limit.

The two large share counters contain 490 and 525 threshold occurrences.
Their threshold lists consist of fourteen terms for each of 35 corner
occurrences, followed in the larger counter by the fifteenth terms.
An inductive list-permutation proof moves these last terms beside their
respective first fourteen terms. It preserves multiplicities, including
corners shared by several centres. Thus it requires no distinct-corner
assumption. The general reordering proof passes Lean.

At the end of round 186, the endpoint counter, all 72 totalizer conditions, and 38 of 98 concrete clause batches had passed. The initial-database truth theorem remains incomplete.
The full 52,489-deduction critical Lean refutation and its interpretation
for actual snakes remain pending. The exact 12,424-clause input is
unchanged. This progress reduces the verification work still required;
it does not establish a new numerical upper bound.

At the end of round 187, 93 of 98 complete clause batches (11904 of 12,424 clauses) have passed Lean. The initial-database truth theorem remains incomplete. The full critical contradiction and the actual-snake floor-fifteen theorem remain pending.

\paragraph{The conditional 233 comparison and weaker share requirements.}
The contributor's six-step argument now supplies a conditional bound 233.
Its saved source and Linux compilation record have been inspected;
the production local certificates and an independent rebuild remain
unfinished. A new deduction reduces the share requirements in that argument.

Keep its original table weights and write their sum as $W$.
There are $m-7$ windows. Suppose each table share value exceeds the physical
five-centre share by at most $q$. The omitted six internal counts and two
endpoints contribute at least $14+8=22$, while every physical share is
counted at most five times. Consequently
\[
W\le120(T-22)+5\sum_i\sigma_{12}(x_i)+q(m-7).
\]
The exact identity in Section 2 then gives
\[
W+4200m+5L\le1098840+q(m-7).
\]
The original potential supplies $520(m-7)\le W+1023$. Setting $q=4$ yields
\[
\boxed{4716m+5L\le1103475\quad\Longrightarrow\quad m\le233.}
\]
Indeed, $m=234$ makes the left side at least $1103544$, exceeding the
right side by 69. This uses the stronger boundary contribution to gain
960 units, permitting weaker local shares without changing the potential.

Five units of allowance would suffice if a smaller potential range of
784 were certified. With range 1023, five units instead force $L\le31$
for any remaining $m\ge234$. The earlier domination theorem then forces
every cube vertex to have a selected neighbour. Thus another sufficient
route is to prove the four-unit requirements in this low-loss case and
the five-unit requirements generally. The stronger low-loss shares are
still unproved. All these numerical implications pass Lean; their
integration with the contributor's actual-snake argument is pending.

In the saved table, four-unit allowance makes 31 of 490 paid share
obligations follow from the four share-sixteen floors. Two established
share-twelve floors cover another 33 rows, leaving 426. Completing the
critical $012034$ floor fifteen would leave 373. These are native table
counts, with the Lean lookup and coverage proof still pending.

\textbf{6. The construction relation between dimensions.}
The equalities $50=2\cdot26-2$ and $98=2\cdot50-2$ suggest 194 as
a construction target in dimension nine. They give no recurrence
theorem or optimality argument. Cutting a snake at its $q$ edges in
one coordinate gives $q+1$ lower-dimensional path components, whose
lengths sum to $m-q$. This follows simply by partitioning its edges
into crossing edges and edges within the two faces.

For the published representatives, cuts with the fewest crossings give
\[
26=2+13+9+2,\qquad
50=18+24+6+2,\qquad
98=8+25+10+25+26+4.
\]
The 98-edge representative crosses every coordinate at least four
times, so its construction cannot be explained by one or two crossings.
One face contains two 25-edge components, or 52 vertices. This exceeds
the 51 vertices allowed in one seven-dimensional snake, illustrating
why a face must be treated as a collection of paths. These statements
describe the given representatives, without classifying all optima.

An exact construction rule is available. Take three induced paths
$P,Q,R$ of lengths $a,b,c$ in $Q_n$. Require that $P,Q$ share only
their joining endpoint, that $Q,R$ share only their joining endpoint,
and that $P,R$ are disjoint with no edge between them. Placing $P,R$
in one face and $Q$ in the other gives a snake of length $a+b+c+2$.
Within each face the assumptions exclude chords; across the faces,
equality of projected vertices creates exactly the two crossing edges.
This rule and the concrete $18+24+6+2=50$ reconstruction pass Lean.

Of 131 saved 191-edge snakes, 84 have the same 95-edge middle-path
type in a two-crossing decomposition, up to cube symmetries and
reversal. This path is not a segment of the displayed 98-edge snake.
It decomposes as $95=31+47+15+2$; its two outer components in $Q_9$
total 94 edges, giving $191=95+94+2$. These are native finite-data
comparisons. They suggest retaining compatible components and their
endpoints when passing to the next dimension.

Cube vertices have two colours according to Hamming-weight parity,
and every edge changes colour. Hence fixed endpoints force the parity
of a replacement path's length. The alternating-colour deduction
passes Lean. A 97-edge replacement for the 95-edge middle path with
the same endpoints and fixed outer components would give 193 edges.
It still must avoid all unintended intersections with those components.
A 194-edge two-crossing construction would require component lengths
totalling 192, for example $97+95$ across the middle and outer parts.

The fixed 50-edge middle-path search reports no two-crossing lift
of at least 98 edges; this is a restricted solver exclusion without
a Lean certificate. A bounded search with the 95-edge middle path
fixed and both outer paths nontrivial leaves the 192-edge target
unresolved. No longer construction has been found. The restrictions obtained from this replacement investigation are
described below.


\paragraph{What obstructs a replacement in the common family.}
The family has 48, 96 and 192 selected vertices at lengths 47, 95 and
191. Its balanced face sizes explain this observed doubling, without
establishing a general recurrence.

The direct fixed-endpoint 95-to-97 search remains unresolved. A smaller
attempt retains the 31- and 15-edge pieces and replaces the 47-edge
component by a 49-edge path. The solver excludes this even after all
external nine-dimensional restrictions are removed. This exclusion
concerns those fixed pieces and has no Lean contradiction certificate.

The displayed optimal lower-dimensional seeds are also far from the
required compatibility. All 1440 endpoint-preserving placements of
the applicable 97-edge truncation overlap the fixed outer vertices
at least 23 times, across all 84 contexts. At the nested level all
960 placements of the 49-edge truncations have at least eleven
forbidden overlaps. The general inequality
\[
|A\cap B|\le |A\setminus A'|+|B\setminus B'|
\quad\text{when }A'\cap B'=\varnothing
\]
passes Lean: each initial overlap must be removed from at least one
set. Here the designated crossing vertices are omitted from the ground
set. The finite placement counts are native computations.

An exhaustive search excludes a two-edge improvement obtained by replacing
any single interval containing at most 42 internal vertices of the common
95-edge path. A second search covers two separated changed positions,
including reversal of the retained middle section, and also finds
no replacement. Deleting two separated vertices leaves three chains;
their two new connections must have lengths $(2,4)$, $(3,3)$ or $(4,2)$.
These cases were completely enumerated. Together with the adjacent
cases, this means a fixed-endpoint 97-edge replacement must omit at
least three original vertices and add at least five. The complete
Lean interpretation of these finite exclusions remains pending.

A further search requires a 97-edge path with the same endpoints
to retain exactly 93 of the original 96 vertices. All other cube vertices
are available, and no external nine-dimensional restriction is imposed.
The solver reports infeasibility in 28.641 seconds. Combined with the
smaller complete enumerations, this gives the computational conclusion
that at least four original vertices must be omitted. This last step
depends on the solver exclusion, which has no Lean contradiction
certificate.

The corresponding search retaining exactly 92 original vertices
also reports infeasibility. Under these solver results, a fixed-endpoint
97-edge replacement must omit at least five original vertices and add
at least seven new vertices. The numerical restriction is specific to
this path and retains the stated certificate limitation.


Further searches can allow a larger contiguous replacement, changes
in separated parts of the path, or changes to the surrounding components
and crossing vertices. No longer snake or improved global bound has been
obtained.


\paragraph{A general restriction on replacements.}
Let $D$ be the vertices omitted from an induced path and $A$ the vertices
added to a longer induced path with the same endpoints. If the gain is
$\delta$ edges and $|D|=r$, then $|A|=r+\delta$. Writing $C$ for the
common vertices, the internal degree sums are
\[
2r=2e(D)+e(D,C),\qquad
2(r+\delta)=2e(A)+e(A,C).
\]
The common induced edges cancel when the two path lengths are subtracted.
Consequently
\[
e(A)=e(D)+\delta,\qquad e(A,C)=e(D,C)=2(r-e(D)).
\]
Since a union of paths has vertex count minus edge count components,
$D$ and $A$ have the same number of components. Lean checks the finite
edge accounting and the deduction from the displayed degree and size
premises; the full arbitrary-path integration remains separate.

For five omissions and a two-edge gain, the numbers of deleted internal
edges are $4,3,2,1,0$ according as the omissions form one through five
intervals. The seven added vertices then induce $6,5,4,3,2$ edges.
The one-interval case was already excluded. An exhaustive search of
all 16020 two-interval placements, both middle-chain orientations and
all compatible connector lengths excludes that case as well.
For three or more intervals, a degree-and-square relaxation reports
infeasibility. It requires the path degrees, the exact replacement
counts and edge identity, and at most three selected vertices per
coordinate square. These are necessary conditions for every desired
path, so infeasibility excludes the remaining cases.

The round-190 computational restriction was at least six omitted
original vertices and eight added vertices for a fixed-endpoint
97-edge replacement. The native exclusions and solver result still
lack a complete Lean contradiction certificate. They do not establish
optimality of the 95-edge component or the 191-edge snake.

A 120-second overlap-maximization run returns the supplied known
97-edge seed, retaining 52 of the original 96 vertices. No closer path
was found, and optimality was not established. The seed has 30 forbidden
overlaps in the first surrounding configuration, so it does not yield
a nine-dimensional improvement. Across all endpoint-preserving placements
of this displayed seed family, 52 is the largest retained intersection.
The lower restriction of six omitted vertices at that stage and the
explicit example omitting 44 left a substantial gap for the minimum
number of omissions. The lower restriction retains the computational
verification qualifications above. A larger contiguous replacement,
separated changes, and revised crossing vertices remain possible.



\paragraph{Current replacement frontier.}
At every retained vertex $v$, the number of adjacent added vertices
equals the number of adjacent deleted vertices:
\[
|N(v)\cap A|=|N(v)\cap D|.
\]
Both paths give $v$ the same degree, and its retained neighbours are
unchanged. Subtracting those neighbours proves the equality. An added
vertex therefore cannot touch a retained vertex that lost no neighbour.

The cube's two colours alternate along a path. A 95-edge path has
48 vertices of each colour, and a 97-edge path has 49 of each.
Consequently the replacement must add one more vertex than it deletes
in each colour. Lean checks the finite counting and alternating-sequence
arguments; the full graph-to-list interpretation remains separate.

A model allowing all cube vertices now excludes replacements omitting
at most six or at most seven original vertices. It fixes the endpoints,
selects 98 vertices, requires the path degrees and colour counts, and
imposes the exchange identity and all square exclusions. These
conditions hold for every desired replacement. Their infeasibility
therefore gives the computational restriction
\[
|D|\ge8,\qquad |A|\ge10.
\]
The eight-omission query remains unresolved after 120 seconds.
The exclusions have no Lean contradiction certificate. The known
97-edge placement omitting 44 vertices remains the explicit upper
example for this local distance problem, so the remaining range is
$8\le r\le44$, with the lower end subject to that qualification.

A separate query retains the outer components of the first saved
nine-dimensional configuration while allowing the middle path to change
without an omission limit. The 97-edge target remains unresolved after
120 seconds. This supplies neither a construction nor an exclusion.


The lower-dimensional comparison also shows why compatibility matters.
With the displayed optimal 26-edge $Q_6$ path fixed as the middle,
the solver gives a maximum 48-edge two-crossing lift, whereas the
known 50-edge construction uses a shorter 24-edge middle:
\[
2+26+18+2=48<50=18+24+6+2.
\]
This is a restriction for the displayed seed, without a classification
of all optimal 26-edge paths. With the displayed 50-edge middle in
$Q_7$, an 89-edge lift is found; its restricted optimum remains
between 89 and 97 using the earlier exclusion. The 192-edge target
with the displayed 98-edge middle in $Q_8$ remains unresolved after
180 seconds. These computations support searching for compatible
middle and outer components together. They provide no new record.



\paragraph{Changing endpoints and crossing vertices.}
A 191-edge path has endpoints of opposite colours, whereas a
192-edge path has endpoints of the same colour. A one-edge improvement
must therefore change at least one global endpoint.

Let $d$ count omitted old endpoints and $a$ count new endpoints among
the added vertices. The replacement identities become
\[
e(A)+a=e(D)+\delta+d,\qquad e(A,C)+d=e(D,C)+a.
\]
Indeed, the degree sums of the omitted and added vertices are
$2|D|-d$ and $2|A|-a$. Substituting these in the previous edge-count
argument gives the corrections. Their finite counting consequences
pass Lean. An old endpoint retained as internal must acquire one
extra neighbour after accounting for deleted neighbours.

A joint search now permits both endpoints, both faces and all crossings
to change. It computationally excludes a 192-edge path omitting at
most four vertices of the selected 191-edge witness. Thus at least
five omissions and six additions are needed in this local problem.
This solver result lacks a Lean contradiction certificate. The case
with at most six omissions and two crossings remains unresolved.

A different investigation considers replacing any internal vertex
$v_i$ by the other square corner
$v_{i-1}\oplus v_i\oplus v_{i+1}$. A conflict with a retained vertex
forces another substitution; conflicts between two new vertices
forbid that pair of substitutions. A proposed endpoint extension
supplies further forced and forbidden choices. Propagating the forced
choices decides this restricted problem: if no conflict occurs,
leaving every other position unchanged gives a valid assignment.

All 2358 endpoint candidates across the 131 saved snakes are excluded
by this native calculation, even with arbitrarily many simultaneous
square substitutions. The 147 valid closures produced by forcing one
substitution all lead to existing bank classes. These 147 cases do
not enumerate every valid simultaneous substitution set.

For the selected 191-edge witness, the corresponding geometric
exclusion is checked in Lean for both endpoints and all nine extension
directions. The proof permits arbitrary Boolean choices at all 190
internal positions and checks the actual alternate-corner coordinates.
It uses only nonconsecutive vertex separation and separation from the
proposed extension vertex, which are necessary conditions for an induced
path. All 18 candidate contradictions and their combined statement pass.
The complete 131-witness census remains a native computation; this Lean
result certifies the selected witness, without asserting its optimality.


Further construction searches can use larger replacement regions,
reordered components, or successive batches with recomputed corners.
The one-batch exclusion does not settle those possibilities.
No longer snake or improved global upper bound has been obtained.



\paragraph{Larger contiguous replacements.}
The completed census now replaces every block of 2--33 edges in each
of the 131 saved 191-edge snakes by every induced route of the same
length between its endpoints. Every valid result remains in an
existing symmetry class. This gives a stronger consequence than a
single failed move: any finite sequence of these replacements stays
in the bank and therefore cannot unlock an endpoint extension.
The general induction passes Lean; the numerical bank-closure premise
remains a native computation.

The replacement criterion is direct. Its internal vertices must avoid
all retained vertices and all neighbours of retained vertices other
than the two anchors. Its own inducedness controls contacts with the
anchors. This excludes every possible repeat or chord, including for
blocks touching a global endpoint. Lean checks the sufficient criterion
for replacements of unchanged length.

Shortening a block by two edges was also examined for all block
lengths 4--24. The resulting 189-edge paths form 231 symmetry classes:
220 allow no endpoint addition and 11 allow only one. None regrows to
191 through endpoint additions alone. Rerouting a shortened path
before regrowth remains open.

The next search can focus on the old vertices that block a proposed
new endpoint. Every such vertex must lie inside the replaced block.
If their first and last indices are $r,s$, the block needs at least
$s-r+2$ edges. This necessary condition and its underlying retained-vertex
argument pass Lean. It directs larger replacements towards the
192-edge target; it does not prove their existence. Complete native
coverage currently ends at 33 edges, and no global bound has changed.


\paragraph{Targeted growth and the lower-dimensional comparison.}
Fixing the proposed new endpoint allows the search to remove its
blocking vertices first. Repeatedly delete non-anchor vertices with
at most one available neighbour; an internal path vertex needs two
distinct path neighbours. Reject disconnected anchors or insufficient
colour counts. Attached graph branches cannot be used without
revisiting their attachment vertex. If a required section has adjacent
entry and exit vertices, inducedness forces their edge directly.
The key implications pass Lean; finite coverage remains native.

These bounds and selected-vertex solver models exclude every single
replacement of at most 44 edges immediately followed by one endpoint
addition, for all saved 191-edge snakes. This is a one-operation
cutoff. The earlier repeated-move closure still extends only through
33 edges.

A lower-dimensional comparison prevents interpreting these exclusions
as optimality. The known suboptimal 48-edge $Q_7$ lift cannot improve
to 49 by this operation with a block of at most 37 edges, yet an
explicit 43-edge replacement succeeds. Both paths and their retained
positions are checked in Lean; the numerical lower cutoff is native.
Thus the required block size in this example lies between 38 and 43.
Large local obstructions can coexist with a better global construction.
The proposed rerouting search on the shorter189-edge states remains
incomplete. The current dimensional route is described next.


\paragraph{Transferring the dimensional construction.}
The useful patterns are compatibility in $6\to 7$ and multiple
components per face in $7\to 8$. In the 98-edge example, one face
contains two 25-edge paths, hence 52 vertices; the other contains
three paths with 47 vertices. Four crossings join these five
components. For $8\to 9$, two 49-edge paths in one face and 95
compatible vertices in three opposite-face paths would give 194
edges. This remains a construction target.

The two 25-edge components are exchanged by an involution that swaps
two coordinates and flips three others. Identifying paired vertices
reduces the search to a smaller graph. An induced path there lifts
to two disjoint, nonadjacent paths: a repeat or a nonconsecutive
contact would project to a repeat or chord, while the distance
between partners prevents two lifts of a consecutive edge.
The lifting argument and this symmetry's finite properties pass Lean.

A restricted construction using 1536 placements reaches two 42-edge
paths in $Q_8$, giving86 face vertices. The saved $Q_9$ bank instead
provides symmetric47-edge pairs under other involutions. Their
quotients require an additional rule: both endpoints of an edge
that lifts to a square cannot be selected. A Lean separation lemma
supports this version under explicit safe-step hypotheses.

The larger quotient and four-crossing searches remain unresolved.
Known 98- and 191-edge patterns are reconstructed when their face
selections are fixed. The next attempt starts from an observed
47-edge pair and seeks a larger pair while excluding disconnected
cycles as they appear. A larger pair must still be joined to a
compatible opposite-face forest before it yields a longer snake.


\paragraph{Removing disconnected cycles from the construction search.}
For the chosen symmetric 47-edge pair, the revised search permits
both endpoints to move and seeks50 quotient vertices. Its degree
constraints can produce a path together with cycles. Whenever a
cycle appears, at least one of its vertices must be excluded.
The reason is direct: the greatest-indexed cycle vertex on an
induced path would need two distinct neighbours at the same
preceding position. Lean checks this argument and all three
36-vertex cycle sets used in the four-omission calculation.

That calculation is infeasible even with free endpoints. Hence any
desired50-vertex safe quotient path under this involution must omit
at least five seed orbits and add at least seven, corresponding to
ten omitted and fourteen added vertices in the paired forest.
The numerical infeasibility remains a native solver result.
The six-omission retry returned UNKNOWN after 150.062 seconds. The six-omission case remains unresolved.
No longer nine-dimensional snake follows. A larger pair must still
be completed by a compatible opposite-face forest.


\paragraph{Using symmetry to organize the next construction.}
A different involution gives the quotient used in the current endpoint
searches. In coordinates 0 through 6, let $h$ exchange 1 with 4 and
3 with 6, then translate by 126. Add $x\leftrightarrow h(x)$ to the
ordinary $Q_7$ edges. Its $Q_8$ extension has minimum displacement three
and no double-lift restrictions. This quotient has three endpoint
classes, represented by 10, 2 and 0. They are distances zero, one and
two from its triangle vertices. Cube symmetries act transitively in
each class, so three unrestricted endpoint cases cover every quotient
path. The corresponding 50-vertex searches remain unresolved after
45 seconds per case. This is a time limit, without an exclusion.

The checked pair of separated 47-edge $Q_8$ paths remains available.
Its 96 vertices require 97 compatible opposite-face vertices for a
192-edge $Q_9$ snake. That completion remains unresolved.

A broader use of symmetry allows any cube isometry and different
parent shapes. Let $S,T$ be the vertex sets of two induced $Q_7$ paths.
The copies $(S,0)$ and $(f(T),1)$ in $Q_8$ are separated exactly when
\[
S\cap f(T)=\varnothing.
\]
Within each face the paths remain induced. Across faces the Hamming
distance is one plus the distance of the projections, so an unwanted
edge is exactly a coincident projection. No involution condition is
needed for this construction. Involution quotients remain a separate
method, since applying $(y,b)\mapsto(f(y),1-b)$ twice gives $f^2(y)$.

For a fixed parent pair, every possible placement is obtained by
5040 coordinate permutations and 128 translations. For each permutation
$\pi$, counting pairs with $u\oplus\pi(v)=t$ gives exactly the
intersection size at translation $t$. Scanning all 645120 entries
therefore decides disjoint placement for that pair. Normalizing either
parent and reversing their order only changes which isometry is scanned.

Extensions of two fixed 24-vertex prefixes of the published parent
give three classes with 50 vertices and fourteen with 49 vertices.
Every placement of any pair from the three 50-vertex classes overlaps
in at least two vertices. Pairing the published 51-vertex parent with
the fourteen 49-vertex parents also always overlaps: the minimum is
two in twelve cases and three in the remaining two. Consequently,
changing relative symmetry alone cannot produce a 100-vertex paired
face from these saved families. The calculation does not cover all
possible parent shapes.

The generic path-image, layer-separation and representative-change
arguments pass Lean. The complete numerical censuses remain native
computations, without a Lean replay.

The trimming family now has an exact native limit. Two independently
trimmed copies of the published 51-vertex parent give at most 92
separated vertices, over all cube isometries. To see why the search is
complete, let $C$ count shared vertices and $D$ count deleted parent
vertices. Each shared vertex requires a deletion, so $C\le D$.
A retained total of at least 92 forces $D\le10$. Only 1363 isometries
have $C\le10$. For each, all 66 first-parent trims of cost at most ten
are examined; the best second interval is the largest gap avoiding
their collisions. The 89958 resulting cases attain 92 and never exceed it.
The 96-vertex forest found earlier changes the parent shape and is
therefore outside this trimming family.

For a repaired pair with additions, let $A$ count added vertices.
A 100-vertex target obeys $D=A+2$. Together with $C\le D$, this leaves
exactly sixteen starting isometries for repairs with $D\le4$, or fifteen
cases after identifying inverse maps. The native catalogue currently excludes 2 of the fifteen inverse-pair representatives at four total omissions; 13 representatives remain unresolved. The two minimum-intersection cases are among the exclusions. No complete four-omission exclusion is claimed.
The overlap-cover inequality, repair balance and interval bounds for
the two minimum-overlap cases pass Lean. The catalogue outcomes remain
native solver results. Six-omission searches in those two placements
remain unresolved after 60 seconds each.


The next broader query seeks any 49-vertex induced path in the
complement of the published 51-vertex parent. This allowed graph has
77 vertices and 141 edges; the query remains unresolved after 60 seconds.
It imposes no symmetry or shape condition on the second parent.
A witness would give a 100-vertex paired face, still needing 93,
95 or 97 compatible opposite-face vertices for 192, 194 or 196 edges.
Further work must change the parent shapes or the face arrangement,
or use matching steps in the earlier quotient. No longer snake or
improved global numerical bound has been obtained.


\paragraph{Fundamental restrictions on dimensional constructions.}
Let two positive-length parent paths of lengths $p,q\le a(n)$ lie
in opposite faces. Delete $D$ original vertices and add $A$ vertices
outside the original copies. Assume the entire final induced graph is
a path of length $L$ with $k$ crossings between faces. Let $b$ count
deleted original endpoints, $e$ count original within-parent edges
deleted at both ends, and $H$ count final within-face edges touching
an added vertex, each edge once. Then
\[
\boxed{L+k+b+e+H=p+q+2+2A,\qquad L+k\le2a(n)+2+2A.}
\]
The proof is an edge count. Initially there are $p+q$ within-face
edges. Deleting vertices removes $2D-b-e$ of them. Adding the $H$
new within-face edges and $k$ crossings gives all final edges, while
the final vertex count is $p+q+2-D+A=L+1$.

Without additions, the same calculation gives $k=1+D-b-e$.
If $C$ is the original projected intersection, its cross-face edges
form a matching. Every destroyed crossing needs a deletion, so
$C\le D+k$ and $D\ge\lceil(C-1)/2\rceil$.
These restrictions apply to two specified parent paths; an arbitrary
higher-dimensional snake need not have such a representation.

The 50-edge $Q_7$ representative has a two-crossing split and is not
excluded by the resulting bound $k\le4$ for two $Q_6$ parents.
For the saved 98-edge $Q_8$ representative, every split except one
has at least five crossings, violating $98+k\le102$.
The remaining split has four crossings and face sizes 47 and 52.
A $Q_7$ parent has at most 51 vertices. Thus this witness cannot
be obtained by deletions alone from two $Q_7$ parents in any
coordinate split. The concrete coordinate obstruction passes Lean.
This conclusion explains a limitation of extrapolating the earlier
construction through two complete parent paths.

For 196 edges in $Q_9$, deletion alone from two parents of length
at most 98 permits at most one deletion and two crossings.
An internal deletion splits one 98-edge parent into two pieces.
The intact 98-edge parent must join these pieces through its two
endpoints, and those two attachment points must be the only retained
projected intersections. The initial intersection therefore has size
two or three. This is an exact criterion, with a written necessity
and sufficiency proof in the main paper.

If four or more crossings are required, at least one addition is
necessary. With exactly one addition, equality forces two 98-edge
parents, two isolated internal omissions, exactly four crossings and
$b=e=H=0$. The added vertex has no same-face neighbour and is
therefore a global endpoint. Removing it leaves a 195-edge,
three-crossing snake with $C\le5$. At $C=5$, every omitted shared
projection retains its opposite copy. That copy prevents attachment
at an endpoint exposed by a deletion, so the new endpoint must
attach at an original parent endpoint. If both final endpoints must
be original vertices, at least two additions are required instead.
Existence of the one-addition alternative remains unresolved.

Six new Lean modules check 24 theorem statements or dependency reports,
covering the list accounting, endpoint geometry, overlap equality and
concrete coordinate obstruction. The complete interpretation for arbitrary
geometric constructions and the one-deletion classification remain
written proofs. The next theoretical question is whether suitably
placed 98-edge parents satisfy the exact one-deletion criterion, or
produce the necessary 195-edge predecessor. No new witness or global
bound follows from the restrictions alone.


\paragraph{A symmetry acting on the whole snake.}
Choose two coordinates, write a vertex as $(z,u)$, and let
$\tau_\epsilon$ interchange the two square words of parity
$\epsilon$, fixing the other two. For a Boolean selector $f$, define
\[
F_f(z,u)=(z,\tau_\epsilon^{f(z)}(u)).
\]
This is a single involution defined on every cube vertex. For a snake,
collect the base words $z$ having a selected word of parity $\epsilon$.
Then $F_f$ preserves every edge and every nonedge of the entire snake
exactly when $f$ is constant on each component of the cube graph on
those base words. The reason is that a selector change across adjacent
base words reverses equality of the two active square words, changing
an edge into a nonedge or conversely. All other pairs keep their status.
Lean checks this criterion and constructs the full induced image path.

The components can be reflected independently. Reflecting a subset or
its complement gives images related by a global square reflection.
With at most three components, the unchanged image and the individual
component images therefore cover every choice up to cube isometry.
A checked 13-edge $Q_5$ example changes one selected vertex and gives
a whole snake inequivalent under cube isometries, including reversal.

The saved 98-edge $Q_8$ parent has a connected active graph for all
56 coordinate-pair/parity choices. This particular family therefore
cannot change its shape beyond ordinary cube symmetries.
For the 131 saved 191-edge $Q_9$ snakes, all 9432 choices have at most
three components. The 408 single-component images all belong to the
existing bank. Complement symmetry and invariance under cube isometries
then imply closure under arbitrary iterations of these maps. These
finite census premises remain native computations.

A self-symmetry preserving the same finite snake as a set is identity
or reversal. Within this controlled-reflection family, a nontrivial
self-map agrees on the whole snake with a global square reflection.
More varied embeddings are possible, as the checked example shows.
Thus the definition is useful, but this family does not escape the
saved $Q_9$ bank or produce a new $Q_8$ parent from the saved one.

The parent analysis also sharpens the saturated $C=5$ case above.
The 195-edge predecessor has one omission in each parent. Their indices
have equal parity. In the one-addition case with $C=5$, at least one
omission lies exactly three edges from an original parent endpoint,
so its index is 3 or 95. Both indices are odd. The short reason is
that the two exposed ports and the retained opposite copy of the
omitted vertex form a complete two-edge component in the other face.
This last conclusion has a written proof and remains necessary only.

\paragraph{Larger blocks and a complete rigidity result for one parent.}
Replace the square reflection by any involutive cube isometry $h$
on a coordinate block and set $F_f(z,u)=(z,h^{f(z)}(u))$.
Join neighbouring base fibres when their selected block words meet
the same two-element orbit of $h$. The whole snake is preserved
exactly when $f$ is constant on every component of this graph.
Indeed, only different selectors across such fibres can change
equality of block words, which is exactly cross-fibre adjacency.
The graph itself is invariant under the map.

If no selected block word is fixed by $h$, following consecutive
snake edges forces one selector value throughout. The map then agrees
on the whole snake with identity or global $h$. Any block involution
with a fixed word is a product of coordinate swaps, possibly with
both bits of a pair complemented. Unchanged block coordinates can
be put in the base. For the saved 98-edge $Q_8$ parent, all 56
one-pair, 840 two-pair and 3360 three-pair graphs are connected.
Four pairs leave no base coordinates. This exhausts all cases:
every admissible controlled block involution, and every sequence of
them, leaves this parent's cube-isometry class unchanged. The generic
criterion and rigidity pass Lean; classification and complete case
coverage use written arguments and native censuses.

Two simultaneous complemented swaps do change the saved 50-edge
$Q_7$ parent. Using pairs $(0,4),(1,2)$ and base word 40 changes
its endpoint distance from six to two. Lean checks both entire
51-vertex paths and every pair relation. This gives a different
parent shape at the same length. Complete placement and trimming
censuses still limit original/image and image/image pairs to 92
separated vertices, so this change has not improved the construction.

For the saved $Q_9$ bank, all 198072 two-pair cases produce only
existing bank classes. Combined with the preceding one-pair result,
this family remains closed under arbitrary iterations. Three- and
four-pair cases in $Q_9$ remain unexamined. No stronger global bound
or longer snake follows. The next result examines every order of
block isometry on the saved $Q_8$ parent.

\paragraph{Independent isometries in every coordinate slice.}
Assign an arbitrary cube isometry $g_z$ to each slice and set
$F(z,u)=(z,g_z(u))$. For neighbouring slices with selected block
sets $A,B$, put $r=g_z^{-1}g_w$. The whole snake is preserved exactly
when $r$ fixes every word in $A\cap B$ and
\[
r(B\setminus A)\cap A=\varnothing.
\]
The first condition retains actual crossings; the second prevents
new ones. Within a slice the common isometry preserves every pair,
and non-neighbouring slices cannot create edges. Lean checks this
criterion and the full induced image-path construction.

Choose a shared word $c$ across an actual crossing. Every relative
map fixing it is $r(x)=c\oplus\pi(x\oplus c)$. To fix all shared
words, $\pi$ may permute only coordinates with identical columns
in the matrix of shared words translated by $c$. This enumerates
all relative isometries, of every order. Distinct columns force
identity; otherwise the no-new-crossing condition further restricts
the choices. The column equivalence and rigidity pass Lean.

A complete native enumeration examines all 254 nontrivial coordinate
partitions of the saved 98-edge $Q_8$ parent, allowing independent
cube isometries in every slice. Every admissible image stays in its
original cube-isometry class. Root normalization removes one global
isometry; actual crossings connect the slices, and the complete
relative-map lists enumerate every subsequent choice. Maps agreeing
on all selected words of a slice can safely be identified. The
coverage argument received independent source review, but the full
finite enumeration has not been checked in Lean.

This family is invariant under cube isometries, so arbitrary sequences
of these moves, including changes of partition, also cannot change
this saved parent's cube class. This strengthens the earlier
involution result and exhausts higher-order controlled block rotations
for that parent. It does not classify other 98-edge parents.

In $Q_9$, all 44016 single-three-cycle choices on the saved bank
produce only existing classes. The 194 nonzero phase assignments
represent 34 existing classes. Broader $Q_9$ slice transformations
remain unexamined. A separate checked $Q_4$ example shows that a
rotation producing a valid whole snake once can create a chord on
its second application; admissibility must be evaluated on the
current selected set.

Round 205 advances certification of the 254-partition calculation
as described below. A construction beyond this restriction must
change the selected paths within slices or use different parents.
No longer $Q_9$ snake or improved global numerical bound is obtained.

\paragraph{Whole-path classification with two free coordinates.}
Write the ordered source as $((z_i,u_i))$, where $u_i\in Q_2$
contains the two coordinates allowed to change. For each block
parity, form the graph of occupied base words using base-cube
adjacency, and let its component count be $c_0$ or $c_1$.
There are exactly $2^{1+c_0+c_1}$ ordered induced paths with
these same indexed base words. If $r$ parity graphs are nonempty,
the number of classes under global block isometries is
$2^{c_0+c_1-r}$.

The reason is short. Full parity changes at every path step,
so source and target have one constant block-parity offset.
After a global bit flip removes it, each selected correspondence
acts within a two-element diagonal and extends to identity or
diagonal transposition. Neighbouring active base words must choose
the same transposition: unequal choices toggle a selected pair's
adjacency. Conversely, switches constant on components preserve
every edge and nonedge. This gives the first count. The eight
square isometries have stabilizers of size $2^{2-r}$, giving the
second count. The latter concerns only block isometries, with
index order and all base words fixed.

The component switches commute and each is its own inverse.
They therefore give a symmetry group acting on entire ordered
embeddings, with every selected vertex in a component switched
simultaneously. The action need not fix the original vertex set.

Thus connectedness of each nonempty active graph forces one
global block isometry to carry the entire snake to its target.
Native spanning trees establish this connectedness for both
parities of all 28 coordinate pairs of the saved $Q_8$ parent.
Changing only one or two coordinate histories therefore yields
no new class for that parent. Changing only one coordinate is
even more restricted: the parity argument directly forces one
constant bit flip. The three-coordinate case is resolved below;
different parents and changes of length remain available directions.

The classification has a complete written proof and independent
review. Lean accepts its parity and diagonal-extension components
and the general spanning-tree argument; the full assembled
theorem is pending. All 56 source-specific graph certificates also pass after splitting
the finite assertions into blocks of eight positions. The initial
unchunked attempts exceeded the memory limit.

Lean accepts complete coverage for 254 of the 254 explicit finite models, using 625 of the 625 planned data, elimination and cover modules. This statement concerns assignments to the literal domains and compatibility tables.

\paragraph{Why three coordinates are also insufficient here.}
Every parity-preserving selected correspondence in $Q_3$ that
preserves adjacency and nonadjacency extends to a cube isometry. Label the four
antipodal pairs by their even words. Equal-parity distinct words
have different labels; opposite-parity words are adjacent exactly
when their labels differ. The correspondence therefore induces a
partial injective permutation of four labels. Every full permutation
of those labels is realized by one of the 24 parity-preserving
cube isometries.

Normalizing the common parity offset applies this extension in
each slice of two snakes with matching indexed base words. Hence
every change to three coordinate histories already belongs to the
slice-isometry family. The saved parent's native rigidity result
therefore requires at least four changed histories for a new class,
with global cube isometry and reversal allowed. In four coordinates
the extension principle fails even for whole paths: $(0,1,3,7,6)$
and $(0,1,3,7,15)$ are induced paths with matching indexed parity
and endpoint distances two and four. Lean checks both entire
examples. This establishes the extension theorem's dimensional
boundary; it supplies no new $Q_8$ or $Q_9$ construction. Lean checks the three-coordinate parity/antipodal identities, the induced injective correspondence of selected labels, and two complete four-dimensional paths with different endpoint distances.

The finite-model coverage must still be connected to the actual
cube-isometry domains and compatibility tables. Explicit global
isometries for all 321 surviving leaves are saved and checked
natively on all 99 vertices; their Lean replay is pending. These
proof gaps remain separate from the component classification.
No new numerical bound or longer $Q_9$ construction is claimed.

\paragraph{Arbitrary histories: the current construction frontier.}
The restriction above can be strengthened for the saved parent.
Write a target as $((z_i,u_i))$, retaining its indexed base words
$z_i$ and allowing every possible free-coordinate word $u_i$.
The full snake condition, for every $i<j$, is
\[
d(z_i,z_j)+d(u_i,u_j)
\begin{cases}
=1,&j=i+1,\\
\ge2,&j\ge i+2.
\end{cases}
\]
This follows because Hamming distance adds across disjoint coordinate
blocks. The second line excludes both repeated vertices and unwanted
edges. Consecutive base edges force equal block words, so merging
their indices into intervals gives one variable per interval without
losing any target. Every indexed pair condition remains in the model,
including pairs whose indices merge into one variable.

Global translation and coordinate permutation normalize the initial
block word to zero and the first block transition to the first
coordinate. Interval parities then alternate. Each variable ranges
over all words of its required parity, without a slice-isometry
restriction. All 70 choices of four free coordinates exhaust after
excluding six explicitly witnessed isometric assignments each.
Thus the native calculation requires at least five changed coordinate
histories for a different class of this parent, at the same length
and index matching. It supplies no global length bound.

Lean checks the general pair criterion, its equivalence with the
usual distinctness and adjacency requirements, and interval
substitution soundness. The concrete tables, normalization, search
coverage and isometry witnesses still require Lean certification.
The 70 exclusions retain native-enumeration status.

The five-coordinate calculation now exhausts all 56 choices after
reducing the remaining coordinate permutations. With the initial
word and first transition normalized, permutations of the other
four coordinates form a group of 24 maps. Applying one map to
every interval word acts on the whole snake and preserves every
constraint. Each finite orbit has a lexicographically least tuple,
so retaining such tuples preserves every orbit.

The pruning rule has a short proof. If every earlier domain is
fixed pointwise by a permutation $\pi$, a current value
$u>\pi(u)$ would be the first difference and violate the chosen
order. It can be deleted. The scan must stop as soon as a retained
value satisfies $u<\pi(u)$, since that choice already makes the
comparison strict and leaves later values unrestricted by this
comparison. Lean checks this rule and the finite-minimum argument.

All 56 models exhaust after excluding the witnessed source-isometric
assignments. A different class of this saved parent therefore requires
at least \emph{six} changed coordinate histories at fixed length and
index matching, at the native-enumeration evidence level. The complete
Lean certificate for the concrete searches remains pending. Of the 28
six-coordinate choices, one now exhausts and 27 remain unresolved.
No new parent is found.

There is also a short checked obstruction to reversing this whole
parent. Its edges $(1,3)$ and $(15,13)$ have equal bitwise differences.
Their reversal images would be $(70,68)$ and $(85,21)$, with different
coordinate directions. A cube isometry preserves equality of edge
differences, so it cannot perform that reversal. Lean checks the
argument and the relevant indices in the entire saved source.

The continuation argument preserves progress across search limits.
For a selected variable with domain $D$ and value $a\in D$, the two
domains $\{a\}$ and $D\setminus\{a\}$ partition its possibilities.
On interruption, retain the unfinished descendants and every ancestor's
unexamined remaining values. Soundly exhausted branches contain no
solution. Induction therefore preserves coverage of all unfinished
assignments. Replacing a covered region by its covering descendants
preserves global coverage; an empty frontier excludes solutions only
when the propagation and exhaustion steps are sound. Lean checks
these generic implications, with six dependency reports.

For this source, retaining coordinate histories 1 and 4 now forces
the original cube-isometry class, at the native-enumeration evidence
level. The search resolves 52 interval variables in 51 committed
steps and 1,713,998 nodes. An independent audit finds the complete
receipt chain consistent. The actual native traversal and concrete
model are still outside the Lean coverage proof.

This excludes one particular pair of retained histories. It does not
raise the general requirement from six changed histories to seven.
The pairs $(0,4)$ and $(3,4)$ still have saved unfinished regions.

A stronger geometric restriction arises from the two retained
coordinates. Split the path into constant-free-word intervals.
For distinct intervals of the same free-word parity, equal free
words force every base pair to have distance two. Otherwise a
vertex repeats or a chord appears. In a square each vertex has
only one opposite corner, so both intervals must visit singleton
base vertices at opposite corners.

For one parity, let $M$ count intervals visiting several base
vertices and $N_z$ those visiting only $z$. Choosing all multi-visit
intervals and the larger singleton family from each opposite pair
therefore gives distinct free words, proving
\[
M+\max(N_{00},N_{11})+\max(N_{01},N_{10})\le 2^{r-1}.
\]
For retained coordinates $(0,5)$ in the saved parent, the even
counts are $M=12$ and $(N_{00},N_{01},N_{10},N_{11})=(5,4,10,10)$.
The selected 32 intervals use every even word of $Q_6$.
Thus each of the five singleton intervals at $00$ shares a word
with one at $11$, and each of the four at $01$ shares with one at
$10$. Exactly nine words are shared by pairs of intervals.

In a branch, every even word must have a possible carrier in the
32 selected intervals. A sole possible carrier is forced. Each
outside singleton must also retain a word available to one of its
opposite-corner partners. These deductions strengthen the search
without prescribing which partner is chosen.

Three Lean modules pass with thirteen reports: the square geometry,
generic capacity and support implications, and extraction of the
81 intervals and counts directly from the literal source. The full
model assembly and native exclusion proof remain pending.
The strengthened $(0,5)$ search remains open, as do $(0,4)$ and
$(3,4)$. The other 24 unresolved pairs retain their initial limited
searches. No new parent or global bound has been obtained.
A successful different parent would still need a separate $Q_9$
joining construction.

\textbf{7. What is proved, and what remains.}
The local count exclusions, new boundary counts, physical loss identity,
share packing and conditional arithmetic are Lean-checked. The actual-cube
empty-square argument also proves that an undominated vertex contributes
at least 96 to $L$: its missing-pair graph is connected, with at least
eight edges, or consists of two complete components, which have even more.
Thus $L\le95$ forces domination. The earlier $L\le3$ reduction and sparse
positive-remainder restrictions are checked, although they are unnecessary
for the stronger conditional comparison above.

Round 178 added five accepted Lean modules with twelve theorem dependency
reports, including the actual-snake reduction under explicit local-window
and potential premises. The full finite-graph certificate has not been
checked in Lean.

Both the retained bound 234 and the contributor bound 233 remain conditional. Required local share
floors, retained count restrictions, reference exclusions and the complete
graph-to-snake certificate must still be established. In particular, the
retained model includes a hypothetical six-count floor 25 on $0123104$
under its specified count conditions. A finite graph computation cannot
prove that every snake traces that graph. The older share-eighteen Lean replay
still checks only 33,111 of 145,215 deductions and establishes no final
contradiction. The new floor-fourteen native replays are separate
certificates and do not complete that older Lean replay. Neither a longer construction nor a proof of 196 was obtained.

The current priority is the construction investigation in Section 6.
The preserved upper-bound work remains:
\begin{enumerate}
\item Integrate the improved six-count boundary and four-unit share allowance
into the contributor comparison. Complete its required local certificates
and the coverage proof for the relaxed share table.
\item Complete the Lean contradiction and geometric interpretation for
the required profile $(15,14,14,14)$. The compact critical certificate
contains 52,489 deductions and 5,342 variables.
\item Discharge the retained count restrictions, reference and context
exclusions, graph coverage and complete finite-edge certificate. These
remain essential assumptions of the conditional upper bound.
\item If an alternative argument is pursued, the uniform floor-fourteen
route still needs to remove its loss-zero and loss-one cases. The separate
count-only route would be completed by proving $s_0\ge14$ for a hypothetical
235-edge snake. Neither alternative is required by the new sufficient
profile once its own premises are proved.
\end{enumerate}
Round 193 adds the contiguous-replacement criterion and the general closure argument, with the finite census described above. Round 192 adds endpoint corrections and a geometric Lean exclusion of
coordinated square substitutions followed by extension for one selected
191-edge witness. The full-bank exclusion remains native. The manuscript-wide source review
completed in round 200 corrects a historical merger claim, distinguishes
the two symmetry quotients, and clarifies finite-walk boundary terms.
It is separate from the pending certificate and proof rebuilds.
No unconditional numerical bound or longer construction has been added.

\smallskip
{\small Full proofs and evidence are in the companion manuscript
\texttt{snake9\_research.tex} and research records 175--209.}
\end{document}
